\chapter{Introduction}

\section{Motivation and related work}

Graph colouring is a fundamental problem in computer science. We are given a graph and we would like to colour the vertices using as few colours overall as possible such that no two adjacent vertices get the same colour. Deciding if a graph can be coloured using $k=2$ colours is easy, as it is equivalent to deciding if the graph is bipartite. However, for \(k\ge 3\), it has been known for a long time that deciding whether \(k\) colours suffice is NP-hard \cite{MR4679190-Karp72}.

Contrasting this worst-case hardness result, there has been a lot of work on average-case settings, see for example \cite{MR1344544-Blum95}. In these settings, the input contains randomness, a colouring is guaranteed to exist, and we aim to recover one with high probability. For the fully random case -- when the initial graph is \Erdos-\Renyi and a $k$-colouring is randomly \textit{planted} in it by assigning colours independently at random and removing edges between the colour classes -- \cite{MR1484153-Alon97} provided an efficient recovery algorithm.

The key question then is what happens between worst-case and average-case -- how much can the full-randomness assumptions of \cite{MR1484153-Alon97} be relaxed? Towards this goal, \cite{MR3536556-David16} looked at semi-random models where either the initial graph or the planted colouring is random. Under the crucial assumption that the initial graph is always a two-sided expander -- the second largest and smallest eigenvalues of its normalised adjacency matrix are bounded away from \(1\) in absolute value -- the authors provided a coarse-grained characterisation of when there is an efficient recovery algorithm and when it is NP-hard to recover a full or almost full colouring. Then, \cite{kumar2018finding} further derandomised this result by establishing a variance condition on the degrees between colour classes.

The algorithms in all the aforementioned papers are based on the observation that there is an encoding of the (planted) colouring that is roughly a linear combination of the eigenvectors of the adjacency matrix corresponding to very negative eigenvalues. Therefore, when the adjacency matrix has only a few very negative eigenvalues, an approximate colouring can be decoded after subspace enumeration \cite{kumar2018finding}.

The case of \textit{one}-sided expansion -- when only the second largest eigenvalue is guaranteed to be bounded away from \(1\) -- is therefore much more challenging as it \textit{seemingly} excludes the aforementioned idea, and it is also less studied. For this setting, \cite{bafna-hsieh-kothari} recently showed that for $k=3$ specifically, a small linear-sized independent set can be recovered -- very far away from a full colouring. Their algorithm is based on showing within the sum-of-squares proof system \cite{MR3727623} that under one-sided expansion -- unless the graph is almost-bipartite -- any two $3$-colourings must be sufficiently similar.

For completeness, note that graph colouring can be studied from other perspectives as well. For example, given a $3$-colourable graph, instead of finding assumptions under which a $3$-colouring can be efficiently recovered, \cite{MR4764817-24} provides a more general algorithm that uses $\tilde{O}(n^{0.19747})$ colours.

\section{Results}

In this work, we provide an almost full characterisation of the recoverability of $k$-colourings in \textit{one}-sided expanders.

Key to our results is the concept of the \textit{model matrix} of a graph, which contains the average degrees between the colour classes: the $(i, j)$-entry gives the average number of neighbours with colour $j$ for a vertex with colour $i$.

In \cref{thm:kalg}, we provide an algorithm that given a $k$-colourable one-sided expander graph whose model matrix has no repeated rows, outputs an almost full \kcl. The more general \cref{thm:kalg-partial} recovers a \textit{partial} colouring of size dependent on the relative sizes of the colour classes corresponding to repeated rows in the model matrix.

As corollaries, we show in \cref{thm:2alg} that the two sides of an \textit{almost} bipartite one-sided expander graph can be almost fully recovered. Then, for $k=3$ we prove in \cref{thm:3alg} that when the colour classes are not very unbalanced, an almost full colouring can be recovered. Then, we show in \cref{thm: random-planting-full} that for any $k=O(1)$, when a random colouring is planted in a one-sided expander graph, a \textit{full} colouring can be recovered with high probability.

As a side result, we show in \cref{cor:independentsetrecovery} that in a one-sided expander that contains an independent set with almost half of the vertices, an independent set in size close to that can be recovered. Both \cref{thm:3alg} and \cref{cor:independentsetrecovery} improve on the corresponding results of \cite{bafna-hsieh-kothari} that can only recover a small linear-sized independent set.

Finally, in \cref{sec: hardness} we provide hardness results corresponding to most of our aforementioned results, showing that the assumptions on the model matrix, one-sided expansion, and balancedness are indeed necessary.

\section{Techniques}

Let $\lambda_1, \ldots, \lambda_k$ and $v_1, \ldots, v_k$ represent the $k$ eigenvalues and eigenvectors respectively of the model matrix $M$ associated with the colour classes of the $k$-colourable graph. We then ``lift'' these $k$-dimensional eigenvectors $v_1, \ldots, v_k$ to define their corresponding $n$-dimensional vectors $u_1, \ldots, u_k$: if a vertex has colour $c$, its entry in the vector $u_i$ is the $c$-th component of $v_i$. Then, the core idea behind \cref{thm:kalg} -- which is the basis of most of our results -- is the following:

\begin{enumerate}
    \item In \cref{lem: eigenvectorsdistinguish} we show that if all rows of $M$ are (sufficiently) distinct, then for any two indices $1\leq x<y\leq k$, there is an eigenvector $v_i$ corresponding to a \textit{non-zero} eigenvalue that distinguishes them (i.e. $v_i[x]\neq v_i[y]$). This means that if we can (approximately) recover the vectors $u_i$ corresponding to eigenvectors of $M$ with non-zero eigenvalues, we can (approximately) recover the $k$-colouring by a simple clustering algorithm.
    \item To be able to recover the aforementioned vectors $u_i$, we first show in \cref{lemma: expansionimpliesvariance} that they are \textit{almost} eigenvectors of the adjacency matrix $A$ of the graph with eigenvalues $\lambda_i$ bounded away from zero (i.e. $Au_i\approx \lambda_i u_i$). This result is based on a surprising connection: if the graph is a one-sided expander, then the variance of the degrees between any two colour classes is small -- intuitively, the adjacency matrix is close to a blow-up of the model matrix.
    \item We recover these vectors $u_i$ approximately by subspace enumeration in the eigenspace of the adjacency matrix bounded away from zero. For this procedure to be efficient, we need that the adjacency matrix has only a few very negative eigenvalues -- this is not guaranteed by assumption, as we only make an assumption about the second largest eigenvalue. In \cref{cor:spectral_small_to_small} we show a surprising more general result that guarantees it nevertheless: for any symmetric matrix with non-negative entries, small top threshold rank implies small bottom threshold rank.
\end{enumerate}


To deduce \cref{thm:3alg} from \cref{thm:kalg}, we observe in \cref{lemma: 3colouringfixedaveragedegree} that for $k=3$, given the colour class sizes, the model matrix is fixed due to regularity, and unless one colour class contains almost half of the vertices, it will have no repeated rows. 

The proof for approximate recovery in the random planted setting (\cref{thm: random-planting-partial}) is similar to the proof of \cref{thm:kalg}, but note that while the original graph is a one-sided expander, the resulting graph after planting might no longer be expanding. Hence, to show low variance of the degrees between colour classes, instead of \cref{lemma: expansionimpliesvariance} we exploit full randomness using concentration bounds. Efficient subspace enumeration is guaranteed by \cref{lemma: threshold_rank_kept} that shows that low threshold rank is preserved after planting. To convert an approximate colouring into a full colouring (\cref{thm: random-planting-full}) we adapt \cite{MR3536556-David16} by showing that the expander mixing lemma and small-set vertex expansion -- on which their algorithms rely -- essentially hold for one-sided expansion as well (\cref{lemma: emlonesided}, \cref{lemma: ssveonesided}).

The proof of \cref{cor:independentsetrecovery} relies on the observation that if an independent set is nearly half the size of the graph, regularity ensures most edges go between the independent set and the rest of the graph. Consequently, the $\pm 1$ vector representing this partition is approximately an eigenvector of the adjacency matrix with a very negative eigenvalue.

Finally, the hardness results in \cref{sec: hardness} build on the known \cref{thm:propA1} which shows that without any expansion assumption, it is NP-hard (under the Unique Games Conjecture) to recover even a linear-sized independent set in an almost-bipartite graph.

\section{An introductory example}

To familiarise the reader with the topic, before diving into our results we first provide a solution to a simpler question that was ``left as an exercise to the reader'' in \cite{MR3536556-David16}:

\begin{lemma}[Linear degree random planting full recovery]
    \label[lemma]{lemma: lineardegree}
    Let $H$ be an $n$-vertex graph with minimum degree $\Omega(n)$. Let $G$ be the graph obtained by randomly planting a $3$-colouring in $H$.
    Then, there exists a polynomial-time algorithm that \whp recovers a full $3$-colouring of $G$.
\end{lemma}

The intuition behind \cref{lemma: lineardegree} is that more edges provide more constraints, enabling the recovery of vertex colours through a process of elimination. The linear minimum degree assumption, however, is very strong. Our random planting full recovery result in \cref{thm: random-planting-full} only requires the degrees to be at least a large enough constant, but it assumes $H$ to be a one-sided expander. \cref{thm: random-planting-full} uses the vertex expansion lower bound \cref{lemma: ssveonesided} to ensure that moderately sized vertex sets have a large enough neighbourhood, resembling linear minimum degree.

\begin{proof}[of \cref{lemma: lineardegree}]
    
First we provide the algorithm, then we show that it works with high probability.

\begin{algorithm}[H]
    \caption{Linear Minimum Degree Recovery}
    \label{alg:linear_min_deg}
    \DontPrintSemicolon
    Select $S \subseteq V$ uniformly at random such that $|S| = \Theta(\log{n})$.\;
    Assert that each vertex $x \notin S$ has at least one neighbour in $S$.\;
    \For{each of the $3^{|S|}$ ways of assigning colours to the vertices in $S$}{
        \For{each vertex $y \notin S$}{
            Exclude from $y$'s possible colours any colour assigned to its neighbours in $S$.\;
        }
        Check if the graph is colourable using at most two remaining colours for each vertex via 2-SAT.\;
    }
\end{algorithm}


Let $d_H(x)$ and $d_G(x)$ be the degrees of vertex $x$ in $H$ and $G$ respectively. As each edge is deleted with probability $1/3$, by linearity of expectation we have $\E[d_G(x)] = \frac{2}{3}d_H(x)$ for any $x \in V(G)$. Then, by the Chernoff bound, since $d_G(x)$ is a sum of $\Omega(n)$ independent Bernoulli random variables, with high probability we get $d_G(x) \geq \frac{d_H(x)}{3}=\Omega(n)$ for all $x \in V(G)$.

Since the algorithm goes through all possible ways of assigning a colour to the vertices in $S$, at least one of these will align with the planted colouring, i.e., it will be possible to still $3$-colour the remaining vertices. It only remains to show that, \whp, the assertion in Line 2 of \cref{alg:linear_min_deg} will hold, so we can indeed use 2-SAT:

\begin{align*}
    \P\left[\exists x\;|\;N_G(x)\cap S=\emptyset\right]&\leq \sum_{x\in V(G)} \left(1-\frac{d_G(x)}{n}\right)^{\Abs{S}} \quad\text{by union bound}\\
    &\leq \sum_{x\in V(G)} e^{-\frac{\Abs{S} d_G(x)}{n}} \quad\text{using $1-x \leq e^{-x}$}\\
    &\leq \sum_{x\in V(G)} e^{-2\log{n}} \quad\text{as $d_G(x) = \Omega(n)$ and $\Abs{S}\geq O(\log{n})$}\\
    &= n^{-1}\;.
\end{align*}

The remaining colouring problem, where each vertex has at most two possible colours, can be efficiently solved using 2-SAT. The formulation involves three types of clauses:

\begin{itemize}
    \item Each vertex $x$ with possible colours $a$ and $b$ must be assigned a colour: $x_a \lor x_b$.
    \item Each vertex $x$ must be assigned a unique colour: $\neg x_a \lor \neg x_b$.
    \item Adjacent vertices $x$ and $y$ must not share the same colour $c$: $\neg x_c \lor \neg y_c$.
\end{itemize}

Here, $x_c$ is a boolean variable indicating whether vertex $x$ is assigned colour $c$.

The runtime of the algorithm is polynomial, as the outer for loop goes through $3^{O(\log{n})}=n^{O(1)}$ combinations and 2-SAT is solvable in polynomial time.

\end{proof}
